\section{Introduction}
\label{sec:introduction}

The YOSO model, for \emph{You Only Speak Once}, was proposed to study
protocols built from parties that send one message and then have no
further role~\cite{GentryEtAl2021YOSOFull}.  Its motivating protection
against adaptive compromise is simple.  Prepare the message, erase all
other state, and send.  If speaking reveals where the party can be
attacked, corruption then arrives too late.  The message has already
been sent, and the secrets used to compute it are gone.

Do we really need the erasure?  Reliable deletion may be unavailable,
for instance when a cloud provider retains snapshots or state must be
kept for later inspection.  Suppose that every speaker can be corrupted
immediately afterward, revealing its full retained state.  At first
the answer seems clear.  After everyone has spoken, the adversary knows
everything.  But not everyone need to speak.  An initially honest
party assigned only to help with the computation can remain silent
forever and keep its secrets.
We therefore call the model \emph{You Atmost Speak Once} (YASO).
YOSO already permits silence.  Without erasure, the decision to leave
some roles silent becomes part of the protocol's protection.

\paragraph{A theoretical model for discovering techniques.}
Our purpose is to isolate this resource and ask what secure computation
can achieve under a severe restriction.  The spirit is similar to that
of the private simultaneous messages model, introduced by Feige, Kilian,
and Naor and developed by Ishai and
Kushilevitz~\cite{FeigeKilianNaor1994MinimalModel,IshaiKushilevitz1997PSM}.
We study YASO in this spirit.  A deliberately idealized model can
suggest techniques with uses beyond the setting in which they were
discovered.

\paragraph{The ledger and corruption model.}
A synchronous public ledger supplies a setup string and records a growing
list of registered public keys and messages.  Everyone sees the same
record.  Each party may post at most once, and subsequent corruption
cannot withdraw or overtake an honest posting.  Parties retain their
secret state and randomness.
We call input providers and output recipients \emph{application parties}.
The \emph{computation parties} perform the auxiliary work on their behalf.

Each computation party is independently corrupt at registration with
probability $p$.  Our construction uses $p<1/3$.  This abstracts a
hidden random assignment of roles to a population of machines, some
already compromised.  We give the adversary these initial corruptions
immediately.  Learning their state helps the adversary without giving
the simulator any application inputs or outputs.  An initially honest
computation party is protected from targeted corruption while silent.
Once it speaks, the adversary may corrupt it, and may do so for every
speaker.

Application roles are often tied to particular users.  The adversary
may choose to corrupt any input provider or output recipient at any
time, even while silent.  It may also choose to leave them honest, in
which case their private data must remain protected.

\paragraph{MPC in a DMZ.}
Consider a service offered by one or more companies with private data
$D$.  A client submits an approved computation $F$ and, upon payment,
receives $F(D)$ privately.  The client reaches the MPC service in a
perimeter network (a DMZ).  Each company keeps the machines holding
its data in the clear behind its firewall.  From there it could provision
workers at different hosting providers and publish their public keys on
the ledger.  A key need not reveal its worker's network address or
hosting provider.  The worker's first message may reveal both.  Even if
it deletes its key and shuts down, its provider may retain a snapshot
or be required to disclose one.  YASO allows the adversary to inspect
the speaker's entire state.

Where might the supply of fresh keys come from?  An existing YASO
computation could take fresh random seeds as secret inputs from the
companies, generate a batch of key pairs, and publish $(\pk_i)_i$
along with an encryption of $(\sk_i)_i$ under the already registered
public key of a provisioning machine.  That machine decrypts behind
the firewall, installs one secret key per fresh worker, and can retire
its receiving key without ever posting under it.  If provisioning a
batch in parallel takes time $T$, and an online round trip takes
time $\tau$, a reserve of $\lceil T/\tau\rceil$ rounds' worth of workers
can cover the delay while the next batch is prepared.  This example
trusts the provisioning machine with the entire batch.  The ledger in
our theorem generates its own registered keys and models placement
through the initial corruption rule above.

\paragraph{The task and construction.}
We give a UC feasibility construction for ongoing MPC in this ledger
model.  Each I/O round applies a public Boolean circuit to the private
state and new secret inputs.  It produces the next state, private outputs
for designated recipients, and one common public output.  Public input
strings and public outputs enter the history that determines future
circuits.

Each round uses an independently sampled garbling.  To carry state across
rounds, prepare the successor before evaluating the current circuit.
The successor's sample publishes a public key and embeds the matching
secret key in its garbling.  The current circuit takes that public key
as an input and encrypts its outgoing state under it.  Anyone can carry
the resulting ciphertext forward, and the successor decrypts it internally.
One round of lookahead in circuit selection suffices.  Applications can
insert identity rounds when their next choice depends on their most
recent public inputs and outputs.  Adaptive garbling accommodates the
incoming state ciphertext and successor key being fixed after a garbling
is published~\cite{HemenwayEtAl2016AdaptiveGarbling}.

The adaptive universal sampler of Hofheinz et
al.~\cite{HofheinzEtAl2016UniversalSamplers} supplies the garbling without
entrusting any participant with the garbler's coins.  From public
parameters $U$ and the same sampling description, everyone derives the
same garbling, public keys, and encrypted packets.  The description fixes
the round circuit and the participants' registered encryption keys, or
\emph{delivery keys}, before a fresh committee releases the nonce that
completes it.  Changing a delivery key gives an independent sample in
the sampler's ideal experiment.

An input provider encrypts its secret string under a fresh input key
from the sample and signs the ciphertext together with its public string
for this round and sample.  Its first valid record before the deadline
fixes both.  What the label committees see are the \emph{public bits}
of the accepted or default ciphertext records, incoming state ciphertext, and
successor key.  For each such wire the sampler shares its two labels
among separate committees, encrypting each share under its member's
delivery key.  Honest members of the selected committee post their
secret keys, allowing everyone to reconstruct the selected label.
Honest members of the other committee stay silent and are never reused.
Corrupting the speakers therefore adds no honest shares of the unused
label.  This matters even for a public key.  If both labels on each
successor-key wire were released, an evaluator could have the outgoing
state encrypted to its own key.  On the selected labels the garbling
decrypts the secret inputs, computes the round, and masks each private
output bit with a pad encrypted to its recipient.

Input encryption uses full non-committing encryption so that the
simulator can explain an honest sender's coins or obtain a key for
interpreting a corrupted provider's submission, depending on when
corruption occurs.  Committee shares and output pads use receiver
non-committing encryption.  Full NCE admits constant ciphertext expansion
for long inputs~\cite{BrakerskiEtAl2020ConstantNCE}.
Section~\ref{sec:overview-construction} follows one I/O round through
the construction.

The universal sampler uses indistinguishability obfuscation in the
random-oracle model.  Each round respects the sampler's size bounds,
while a single setup supports an ongoing computation whose number of
rounds need not be fixed in advance.  The aim is to understand what
silence can protect when every speaker's state is exposed, and to
identify techniques for more practical erasure-free protocols.
